\documentclass[aspectratio=169,10pt]{beamer}
\usetheme[numbering=fraction,progressbar=frametitle]{metropolis}
\usepackage{fontspec}
\setsansfont{Latin Modern Sans}
\usepackage{amsmath,amssymb,mathtools,graphicx,booktabs}
\usepackage{polyglossia}
\setdefaultlanguage{french}
\setbeamertemplate{navigation symbols}{}
\setbeamerfont{frametitle}{size=\large}
\setbeamerfont{normal text}{size=\normalsize}
\setlength{\parskip}{.45em}
\hypersetup{pdfauthor={},pdfsubject={Lecture slides},pdfkeywords={},pdftitle={Calcul algébrique, équations et inéquations}}
\providecommand{\N}{\mathbb{N}}
\providecommand{\Z}{\mathbb{Z}}
\providecommand{\D}{\mathbb{D}}
\providecommand{\Q}{\mathbb{Q}}
\providecommand{\R}{\mathbb{R}}
\providecommand{\C}{\mathbb{C}}
\providecommand{\vect}[1]{\overrightarrow{#1}}
\providecommand{\vi}{\vec{\imath}}
\providecommand{\vj}{\vec{\jmath}}
\providecommand{\vk}{\vec{k}}
\providecommand{\norm}[1]{\left\lVert #1\right\rVert}
\providecommand{\vc}[2]{\begin{pmatrix}#1\\#2\end{pmatrix}}
\providecommand{\deter}[2]{\det\left(#1,#2\right)}
\providecommand{\abs}[1]{\left\lvert #1\right\rvert}
\providecommand{\dd}{\,\mathrm{d}}
\providecommand{\e}{\mathrm{e}}
\providecommand{\Card}{\operatorname{Card}}
\providecommand{\E}{\operatorname{E}}
\providecommand{\V}{\operatorname{V}}
\providecommand{\degre}{^{\circ}}
\providecommand{\un}[1]{\,\mathrm{#1}}
\providecommand{\sci}[2]{#1\times10^{#2}}
\providecommand{\pgcd}{\operatorname{PGCD}}
\providecommand{\ppcm}{\operatorname{PPCM}}
\title{Calcul algébrique, équations et inéquations}
\subtitle{Cours et exemples}
\date{}
\author{}
\begin{document}
\begin{frame}[plain]
\vfill
{\usebeamerfont{title}\inserttitle\par}
\medskip
{\usebeamerfont{subtitle}\insertsubtitle\par}
\medskip
{\usebeamercolor[fg]{progress bar}\rule{\textwidth}{.5pt}}
\vfill
\end{frame}
\begin{frame}{Parenthèses et distributivité}
Après un signe \(+\), supprimer les parenthèses conserve les signes des termes. Après un signe \(-\), chaque terme à l'intérieur change de signe. Cette règle concerne une somme ou une différence de termes.
\end{frame}
\begin{frame}{Parenthèses et distributivité}
Supprimer les parenthèses et réduire : \(A=x+(2+3x)\) ; \(B=x-(2+3x)\) ; \(C=-2x+(-4+5x)\) ; \(D=-2x-(-4+5x)\).
\end{frame}
\begin{frame}{Correction : Parenthèses et distributivité}
\(A=4x+2\) ; \(B=-2x-2\) ; \(C=3x-4\) ; \(D=-7x+4\).
\end{frame}
\begin{frame}{Parenthèses et distributivité}
Pour tous réels \(a,b,c,d\) :
\[b(c+d)=bc+bd,\qquad(a+b)(c+d)=ac+ad+bc+bd.\]
\end{frame}
\begin{frame}{Parenthèses et distributivité}
Développer et réduire :
\[\begin{aligned}
A&=-4(8x+7)+(-8x+5)(5x+9),\\
B&=(6x+9)(-x+7)+(3x+5)(-9x+6),\\
C&=-9(-x-7)-(-x+3)(-8x-8),\\
D&=-5(-6x-7)+(x-7)(5x+8).
\end{aligned}\]
\end{frame}
\begin{frame}{Correction : Parenthèses et distributivité}
On développe chaque produit, puis on regroupe les termes de même degré :
\[\begin{aligned}
A&=-32x-28-40x^2-72x+25x+45=-40x^2-79x+17,\\
B&=-6x^2+33x+63-27x^2-27x+30=-33x^2+6x+93,\\
C&=9x+63-(8x^2-16x-24)=-8x^2+25x+87,\\
D&=30x+35+5x^2-27x-56=5x^2+3x-21.
\end{aligned}\]
\end{frame}
\begin{frame}{Factoriser avec un facteur commun}
\(ka+kb=k(a+b)\) et \(ka-kb=k(a-b)\). Factoriser transforme une somme en un produit ; développer fait l'opération inverse.
\end{frame}
\begin{frame}{Factoriser avec un facteur commun}
Factoriser :
\[\begin{aligned}
A&=15x^2-21x,\\B&=3(2+3x)-(5+2x)(2+3x),\\
C&=(4x-1)(x+6)+(4x-1),\\D&=(1-6x)^2-(1-6x)(2+5x).
\end{aligned}\]
\end{frame}
\begin{frame}{Correction : Factoriser avec un facteur commun}
\(A=3x(5x-7)\).

\(B=(2+3x)[3-(5+2x)]=(2+3x)(-2-2x)=-2(2+3x)(1+x)\).

\(C=(4x-1)[(x+6)+1]=(4x-1)(x+7)\).

\(D=(1-6x)[(1-6x)-(2+5x)]=(1-6x)(-1-11x)=(6x-1)(11x+1)\).
Vérifier les facteurs en redéveloppant.
\end{frame}
\begin{frame}{Factoriser avec un facteur commun}
Choisir une forme utile de \(f(x)=(x-2)(x+3)\) pour calculer \(f(2)\), \(f(0)\) et chercher ses zéros.
\end{frame}
\begin{frame}{Correction : Factoriser avec un facteur commun}
La forme factorisée donne immédiatement \(f(2)=0\) et les zéros \(2,-3\). En \(0\), \(f(0)=(-2)\times3=-6\) ; la forme développée \(x^2+x-6\) donne aussi ce résultat. Le choix dépend de la question.
\end{frame}
\begin{frame}{Identités remarquables : développer}
\[(a+b)^2=a^2+2ab+b^2,\qquad(a-b)^2=a^2-2ab+b^2,\]
\[(a+b)(a-b)=a^2-b^2.\]
\end{frame}
\begin{frame}{Identités remarquables : développer}
Développer \((a+b)(a+b)\) : \(a^2+ab+ba+b^2=a^2+2ab+b^2\). Remplacer \(b\) par \(-b\) donne la deuxième identité. Enfin \((a+b)(a-b)=a^2-ab+ab-b^2=a^2-b^2\).
\end{frame}
\begin{frame}{Identités remarquables : développer}
Développer \((x+3)^2\), \((3x-4)^2\), \((x-3)(x+3)\), \((2x+3)(2x-3)\), \((4-3x)(3x+4)\).
\end{frame}
\begin{frame}{Correction : Identités remarquables : développer}
\(x^2+6x+9\) ; \(9x^2-24x+16\) ; \(x^2-9\) ; \(4x^2-9\) ; \(16-9x^2\). Dans le dernier produit, réordonner \(3x+4\) en \(4+3x\).
\end{frame}
\begin{frame}{Identités remarquables : développer}
Développer et réduire \(A=(2x-3)^2+(x+5)(3-x)\) et \(B=2(x+3)+(2x+3)(2x-3)\).
\end{frame}
\begin{frame}{Correction : Identités remarquables : développer}
\(A=4x^2-12x+9-x^2-2x+15=3x^2-14x+24\). \(B=2x+6+4x^2-9=4x^2+2x-3\).
\end{frame}
\begin{frame}{Identités remarquables : factoriser}
Repérer les deux carrés ; pour les formes à trois termes, vérifier le double produit \(2ab\). Pour une différence de deux carrés, utiliser \(a^2-b^2=(a-b)(a+b)\).
\end{frame}
\begin{frame}{Identités remarquables : factoriser}
Factoriser :
\[A=x^2-2x+1,\quad B=25+16x^2-40x,\quad C=4x^2+12x+9,\]
\[D=1-49x^2,\quad E=9x^2-4,\quad F=(2x+3)^2-64.\]
\end{frame}
\begin{frame}{Correction : Identités remarquables : factoriser}
\(A=(x-1)^2\). \(B=(4x-5)^2\). \(C=(2x+3)^2\).

\(D=(1-7x)(1+7x)\). \(E=(3x-2)(3x+2)\).

\(F=[(2x+3)-8][(2x+3)+8]=(2x-5)(2x+11)\).
Pour les trois premiers, vérifier le terme intermédiaire ; pour \(F\), traiter \(2x+3\) comme un seul bloc.
\end{frame}
\begin{frame}{Identités remarquables : factoriser}
Les égalités \((a+b)^2=a^2+b^2\) et \(a^2+b^2=(a-b)(a+b)\) sont-elles vraies pour tous réels ?
\end{frame}
\begin{frame}{Correction : Identités remarquables : factoriser}
Non : pour \(a=b=1\), la première donnerait \(4=2\) et la seconde \(2=0\). Le terme manquant dans la première est \(2ab\) ; le produit de la seconde vaut \(a^2-b^2\).
\end{frame}
\begin{frame}{Puissances entières relatives}
Pour \(n\geq1\), \(a^n\) est le produit de \(n\) facteurs égaux à \(a\). Pour \(a\ne0\), \(a^0=1\) et \(a^{-n}=1/a^n\). Dans ce cours, \(0^0\) n'est pas défini.
Pour \(a,b\ne0\) et des entiers relatifs \(m,n\) :
\[a^ma^n=a^{m+n},\quad\dfrac{a^m}{a^n}=a^{m-n},\quad(a^m)^n=a^{mn},\]
\[(ab)^n=a^nb^n,\qquad\left(\dfrac ab\right)^n=\dfrac{a^n}{b^n}.\]
Pour des exposants strictement positifs, certaines bases peuvent être nulles si les expressions restent définies.
\end{frame}
\begin{frame}{Puissances entières relatives}
Simplifier \(A=2^4\times2^3\), \(B=(-5)^{12}/(-5)^3\), \(C=(7^3)^6\times7^{-5}\), \(D=7^{24}\times3^2\times2^3\). Calculer aussi \((-3)^5\), \(5^1\) et \(99999^0\).
\end{frame}
\begin{frame}{Correction : Puissances entières relatives}
\(A=2^7=128\). \(B=(-5)^9=-1953125\). \(C=7^{18-5}=7^{13}\). \(D=72\times7^{24}\) : on ne peut pas additionner les exposants de bases différentes. Les trois dernières valeurs sont \(-243\), \(5\), \(1\).
\end{frame}
\begin{frame}{Puissances entières relatives}
La définition comme produit répété ne vaut pas pour un exposant négatif. \(-3^2=-9\), tandis que \((-3)^2=9\).
\end{frame}
\begin{frame}{Puissances entières relatives}
Résoudre mentalement \(2^n>100\) pour le plus petit entier naturel \(n\). Écrire un programme qui le trouve.
\end{frame}
\begin{frame}{Correction : Puissances entières relatives}
\(2^6=64\leq100<128=2^7\), donc \(n=7\). Initialiser \(n=0\), \(p=1\) ; tant que \(p\leq100\), remplacer \(p\) par \(2p\) et \(n\) par \(n+1\). Le programme est donné en fin de séquence.
\end{frame}
\begin{frame}{Fractions et expressions fractionnaires}
Pour les dénominateurs non nuls :
\[\dfrac ac\pm\dfrac bc=\dfrac{a\pm b}{c},\quad\dfrac ac\pm\dfrac bd=\dfrac{ad\pm bc}{cd},\quad\dfrac ac\times\dfrac bd=\dfrac{ab}{cd}.\]
Une fraction \(c/d\) possède un inverse seulement si \(c\ne0\) et \(d\ne0\). Alors
\[\dfrac ab\div\dfrac cd=\dfrac ab\times\dfrac dc=\dfrac{ad}{bc}\quad(b,c,d\ne0).\]
\end{frame}
\begin{frame}{Fractions et expressions fractionnaires}
Calculer et rendre irréductible :
\[A=\left(\dfrac23+\dfrac16\right)\times\dfrac45,\quad B=\dfrac12+\dfrac34\times\dfrac56,\quad C=\dfrac45\div\dfrac23\times\dfrac7{10},\]
\[D=\dfrac{\left(\dfrac25+\dfrac13\right)\times\dfrac34}{\dfrac12-\dfrac16}.\]
\end{frame}
\begin{frame}{Correction : Fractions et expressions fractionnaires}
\(A=\dfrac56\times\dfrac45=\dfrac23\). \(B=\dfrac12+\dfrac58=\dfrac98\). \(C=\dfrac45\times\dfrac32\times\dfrac7{10}=\dfrac{21}{25}\). \(D=\dfrac{11}{15}\times\dfrac34\div\dfrac13=\dfrac{11}{20}\times3=\dfrac{33}{20}\). Dans \(C\), multiplication et division sont effectuées de gauche à droite.
\end{frame}
\begin{frame}{Fractions et expressions fractionnaires}
Simplifier \(\dfrac{x^2-1}{x-1}\) en précisant le domaine. Peut-on ensuite prendre \(x=1\) ?
\end{frame}
\begin{frame}{Correction : Fractions et expressions fractionnaires}
Pour \(x\ne1\), \(\dfrac{(x-1)(x+1)}{x-1}=x+1\). L'expression initiale reste non définie en \(1\) ; la simplification ne lève pas la restriction.
\end{frame}
\begin{frame}{Racines carrées et conditions de validité}
Pour \(a\geq0\), \(\sqrt a\) est l'unique nombre positif ou nul dont le carré vaut \(a\). Ainsi \((\sqrt a)^2=a\). Pour tout \(x\in\R\), \(\sqrt{x^2}=|x|\).
Pour \(a,b\geq0\), \(\sqrt{ab}=\sqrt a\sqrt b\). Si \(b>0\), \(\sqrt{a/b}=\sqrt a/\sqrt b\).
\end{frame}
\begin{frame}{Racines carrées et conditions de validité}
Le nombre \(\sqrt a\sqrt b\) est positif ou nul et son carré vaut \(ab\) ; par unicité de la racine carrée positive, il égale \(\sqrt{ab}\). Le même argument appliqué au quotient établit la deuxième formule, sous \(b>0\). De même \(|x|\geq0\) et \(|x|^2=x^2\), donc \(\sqrt{x^2}=|x|\).
\end{frame}
\begin{frame}{Racines carrées et conditions de validité}
Simplifier \(\sqrt{18}\) ; \(\sqrt2\sqrt8\) ; \(\sqrt{7/25}\) ; \(\sqrt{27/75}\) ; \(\sqrt{(-5)^2}\).
\end{frame}
\begin{frame}{Correction : Racines carrées et conditions de validité}
\(3\sqrt2\) ; \(4\) ; \(\sqrt7/5\) ; \(\sqrt{9/25}=3/5\) ; \(5\).
\end{frame}
\begin{frame}{Racines carrées et conditions de validité}
Il n'existe pas de règle générale \(\sqrt{a+b}=\sqrt a+\sqrt b\). Avec \(a=4,b=9\), on aurait à tort \(\sqrt{13}=5\). Dire que ces expressions sont toujours différentes serait aussi faux, par exemple si \(a=0\).
\end{frame}
\begin{frame}{Racines carrées et conditions de validité}
Simplifier \(2\sqrt{18}-\sqrt{50}+\sqrt8\). Encadrer \(\sqrt{20}\) entre deux entiers consécutifs.
\end{frame}
\begin{frame}{Correction : Racines carrées et conditions de validité}
\(6\sqrt2-5\sqrt2+2\sqrt2=3\sqrt2\). Comme \(16<20<25\), \(4<\sqrt{20}<5\).
\end{frame}
\begin{frame}{Équations du premier degré}
Résoudre une équation, c'est trouver toutes les valeurs de l'inconnue qui vérifient l'égalité. Deux équations sont équivalentes si elles ont les mêmes solutions. Ajouter une même expression ou multiplier par un même nombre non nul conserve l'équivalence.
\end{frame}
\begin{frame}{Équations du premier degré}
Développer si nécessaire, regrouper les termes en \(x\), puis isoler \(x\). Si son coefficient est nul, examiner l'égalité restante : elle peut être toujours vraie ou impossible. Vérifier dans l'équation de départ.
\end{frame}
\begin{frame}{Équations du premier degré}
Résoudre :
\[\begin{aligned}
a)&\ 5x-3=2x+9,& b)&\ 3(2x-1)=2(x+4)+7,\\
c)&\ -2(x+3)=4x-(x-1),&d)&\ \dfrac{x+2}{3}=x-1,\\
e)&\ \dfrac{-5x+5}{8}=-\dfrac34,&f)&\ \dfrac{-3x-1}{9}=-\dfrac74.
\end{aligned}\]
\end{frame}
\begin{frame}{Correction : Équations du premier degré}
a) \(3x=12\), donc \(x=4\). b) \(6x-3=2x+15\), donc \(x=9/2\). c) \(-2x-6=3x+1\), donc \(x=-7/5\). d) \(x+2=3x-3\), donc \(x=5/2\). e) \(-5x+5=-6\), donc \(x=11/5\). f) \(4(-3x-1)=-63\), donc \(-12x=-59\) et \(x=59/12\).
\end{frame}
\begin{frame}{Équations du premier degré}
Résoudre \(2(x+1)=2x+2\), puis \(2(x+1)=2x+3\).
\end{frame}
\begin{frame}{Correction : Équations du premier degré}
La première se réduit à \(2=2\) : \(S=\R\). La seconde donne \(2=3\) : \(S=\varnothing\). On ne divise pas par un coefficient nul.
\end{frame}
\begin{frame}{Produit nul et équations avec un quotient}
\(AB=0\iff A=0\) ou \(B=0\). Pour \(B\ne0\), \(A/B=0\iff A=0\). Pour résoudre \(A(x)/B(x)=k\), imposer d'abord \(B(x)\ne0\), puis résoudre \(A(x)=kB(x)\) et écarter les valeurs interdites.
\end{frame}
\begin{frame}{Produit nul et équations avec un quotient}
Résoudre \((4x+6)(3-7x)=0\) ; \(3x(-7x+5)=0\) ; \(\dfrac{2x+1}{3x-5}=0\).
\end{frame}
\begin{frame}{Correction : Produit nul et équations avec un quotient}
Solutions : \(\{-3/2;3/7\}\) ; \(\{0;5/7\}\) ; \(\{-1/2\}\), avec valeur interdite \(5/3\) dans le dernier cas.
\end{frame}
\begin{frame}{Produit nul et équations avec un quotient}
Factoriser puis résoudre \((1-6x)^2-(1-6x)(2+5x)=0\) et \((4x-1)(x+6)+(4x-1)=0\).
\end{frame}
\begin{frame}{Correction : Produit nul et équations avec un quotient}
Première : \((6x-1)(11x+1)=0\), donc \(S=\{1/6;-1/11\}\). Deuxième : \((4x-1)(x+7)=0\), donc \(S=\{1/4;-7\}\).
\end{frame}
\begin{frame}{Produit nul et équations avec un quotient}
Résoudre \(\dfrac{x+1}{x-2}=3\), puis \(\dfrac{x^2-4}{x-2}=4\).
\end{frame}
\begin{frame}{Correction : Produit nul et équations avec un quotient}
Dans les deux cas, \(x\ne2\). Première : \(x+1=3x-6\), donc \(x=7/2\), autorisé. Deuxième : pour \(x\ne2\), le quotient vaut \(x+2\) ; l'équation impose \(x=2\), interdit. Donc aucune solution.
\end{frame}
\begin{frame}{Équations avec un carré ; isoler une variable}
Dans \(\R\), \(x^2=a\) a deux solutions \(\pm\sqrt a\) si \(a>0\), une solution \(0\) si \(a=0\), et aucune si \(a<0\).
\end{frame}
\begin{frame}{Équations avec un carré ; isoler une variable}
Résoudre \(x^2=0\) ; \(x^2=-8\) ; \(x^2=7\) ; \(25x^2=81\) ; \((x+2)^2=9\).
\end{frame}
\begin{frame}{Correction : Équations avec un carré ; isoler une variable}
Dans l'ordre : \(\{0\}\) ; \(\varnothing\) ; \(\{-\sqrt7;\sqrt7\}\) ; \(\{-9/5;9/5\}\) ; \(\{-5;1\}\). Pour la dernière, \(x+2=-3\) ou \(x+2=3\).
\end{frame}
\begin{frame}{Équations avec un carré ; isoler une variable}
Décrire les opérations effectuées sur les deux membres ; préciser les conditions qui permettent de diviser ou de prendre une racine carrée.
\end{frame}
\begin{frame}{Équations avec un carré ; isoler une variable}
Isoler \(h\) dans \(V=\pi r^2h\) avec \(r>0\) ; isoler \(v\) dans \(E=\dfrac12mv^2\) avec \(m>0\), \(E\geq0\) et \(v\geq0\) ; isoler \(x\) dans \(y=ax+b\) avec \(a\ne0\).
\end{frame}
\begin{frame}{Correction : Équations avec un carré ; isoler une variable}
\(h=V/(\pi r^2)\). Puis \(v^2=2E/m\) ; comme \(v\geq0\), \(v=\sqrt{2E/m}\). Enfin \(x=(y-b)/a\). Sans la condition \(v\geq0\), deux signes seraient possibles pour \(E>0\).
\end{frame}
\begin{frame}{Inégalités et comparaison}
On peut additionner des inégalités de même sens : \(a\leq b\), \(c\leq d\Rightarrow a+c\leq b+d\). Ajouter un même réel conserve le sens. Multiplier ou diviser par un réel strictement positif conserve le sens ; par un réel strictement négatif, le renverse.
Comparer \(A\) et \(B\) revient à étudier le signe de \(A-B\). Si \(A,B>0\), on peut aussi comparer \(A/B\) à \(1\).
\end{frame}
\begin{frame}{Inégalités et comparaison}
1. Si \(x\leq-4\), que dire de \(3x-4\) ? 2. Si \(x>-6\), que dire de \(2x-3\) ? 3. Si \(7\leq x\leq16\) et \(3\leq y\leq12\), encadrer \(8x+7y\). 4. Si \(-10\leq x\leq-6\), encadrer \(2x-6\).
\end{frame}
\begin{frame}{Correction : Inégalités et comparaison}
1. \(3x-4\leq-16\). 2. \(2x-3>-15\). 3. \(56\leq8x\leq128\) et \(21\leq7y\leq84\), donc \(77\leq8x+7y\leq212\). Les bornes sont atteignables. 4. \(-26\leq2x-6\leq-18\).
\end{frame}
\begin{frame}{Inégalités et comparaison}
Résoudre \(3x+1\geq x-7\) ; \(2x-11\leq4x-5\) ; \(6x+19\leq2x-3\).
\end{frame}
\begin{frame}{Correction : Inégalités et comparaison}
1. \(2x\geq-8\), soit \([-4;+\infty[\). 2. \(-6\leq2x\), soit \([-3;+\infty[\). 3. \(4x\leq-22\), soit \(]-\infty;-11/2]\).
\end{frame}
\begin{frame}{Signe d'une expression affine et d'un produit}
Pour \(a\ne0\), \(ax+b\) s'annule en \(x_0=-b/a\) ; son signe à droite est celui de \(a\), à gauche l'opposé. Si \(a=0\), le signe est celui de la constante \(b\). Pour un produit, multiplier les signes des facteurs.
\end{frame}
\begin{frame}{Signe d'une expression affine et d'un produit}
Dresser les tableaux de signes de \(3-x\), \(2x-6\), puis \(P(x)=(2x-6)(x+1)\).
\end{frame}
\begin{frame}{Correction : Signe d'une expression affine et d'un produit}
Les deux expressions affines s'annulent en \(3\) : \(3-x\) est positif à gauche, négatif à droite ; \(2x-6\) a les signes inverses.
  \begin{center}\begin{tabular}{llllllll}
\(x\) & \(-\infty\) &  & \(-1\) &  & \(3\) &  & \(+\infty\) \\ 
\(2x-6\) &  & \(-\) & \(-\) & \(-\) & \(0\) & \(+\) &  \\ 
\(x+1\) &  & \(-\) & \(0\) & \(+\) & \(+\) & \(+\) &  \\ 
\(P(x)\) &  & \(+\) & \(0\) & \(-\) & \(0\) & \(+\) &  \\ 
\end{tabular}\end{center}
\end{frame}
\begin{frame}{Signe d'une expression affine et d'un produit}
Résoudre \((7+3x)(2x+6)>0\).
\end{frame}
\begin{frame}{Correction : Signe d'une expression affine et d'un produit}
Les zéros sont \(-7/3\) et \(-3\), à ranger \(-3<-7/3\). Le produit est positif à l'extérieur de ces racines et négatif entre elles. L'inégalité est stricte : \(S=]-\infty;-3[\cup]-7/3;+\infty[\).
\end{frame}
\begin{frame}{Signe d'un quotient et synthèse algorithmique}
Chercher les valeurs interdites, puis les zéros du numérateur. Ranger toutes les valeurs critiques et appliquer la règle des signes. Sur la ligne du quotient, une double barre indique une valeur interdite.
\end{frame}
\begin{frame}{Signe d'un quotient et synthèse algorithmique}
Dresser le tableau de signes de \(Q(x)=\dfrac{-4x+6}{x-3}\), puis résoudre \(\dfrac{2x-1}{-4x+3}\leq0\).
\end{frame}
\begin{frame}{Correction : Signe d'un quotient et synthèse algorithmique}
Pour \(Q\), zéro \(3/2\), interdit \(3\). Signes : négatif sur \(]-\infty;3/2[\), nul en \(3/2\), positif sur \(]3/2;3[\), non défini en \(3\), négatif sur \(]3;+\infty[\).
Pour la seconde expression, zéro \(1/2\), interdit \(3/4\). Les signes sont \(-,+,-\) sur les trois intervalles. Ainsi \(S=]-\infty;1/2]\cup]3/4;+\infty[\).
\end{frame}
\begin{frame}[fragile]{Python}
\small
\begin{verbatim}
def premiere_puissance(seuil):
    # seuil reel >= 1
    n, p = 0, 1
    while p <= seuil:
        n = n+1
        p = 2*p
    return n, p
\end{verbatim}
\end{frame}
\begin{frame}{Signe d'un quotient et synthèse algorithmique}
Que renvoie \texttt{premiere\_puissance(100)} ? Quel lien reste vrai entre \texttt{p} et \texttt{n} à chaque étape ? Comment tester ensuite le signe d'un nombre en Python ?
\end{frame}
\begin{frame}{Correction : Signe d'un quotient et synthèse algorithmique}
Le résultat est \((7;128)\). L'invariant est \(p=2^n\), vrai à l'initialisation puis conservé. Pour le signe, tester d'abord \texttt{x > 0}, puis \texttt{x < 0}, sinon \(x=0\). Ne pas remplacer l'inégalité large de la boucle par une stricte sans modifier le problème posé.
\end{frame}
\end{document}
